\subsection{Finite Taylor graphs and a real normalized comparison}

The formal opening family has two later uses.  Real stopped surfaces
must have the same action coefficients, and actual flow slices must be
compared with a fixed truncation.  Both uses follow from the same local
augmented estimate.  We give it here, after the inverse on the reference
family has been established.

For this subsection take the separated end list of
\Cref{geo:tangent}.  The real family then has the uniform normal tubular
neighborhood supplied by the family construction.  This is the geometry
used by both later applications.

Write $p=(b,x)$ for the reduced coordinates and
$q=\mathsf P(p)$ for the compact observation values.  The observation
functional on a graph is still denoted by $P$.  On a real normal graph
of height $v$ over $S_b$, put
\[
 \mathcal N_b(v,\lambda)
   =\mathcal Q_b(v)+\mathcal C_b(v)\lambda.
\]
Here $\mathcal Q_b$ is the scalar pullback of the shrinker defect and
$\mathcal C_b(v)$ is its compact variational source map.  Thus this is
the sign convention of \eqref{lin:normalized-equation}.  Put
$\lambda_b=\lambda(b,0)$,
$A_b=D_v\mathcal N_b(0,\lambda_b)$, and
$C_b=\mathcal C_b(0)$.  In particular $A_b$ includes the derivative of
the compact source at the generally noncritical member $S_b$.

For $a$ sufficiently close to $1$, let $L^2_a(S_b)$ use
$e^{-a|Y|^2/4}dA_{S_b}$ and set
\begin{equation}\label{ref:weighted-H2}
 \|v\|_{\mathcal H^2_a}^2
 =\int_{S_b}\left(|\nabla^2v|^2+(1+|Y|)^2|\nabla v|^2
                    +(1+|Y|)^4|v|^2\right)e^{-a|Y|^2/4}dA_{S_b}.
\end{equation}
Corollary~\ref{lin:shifted-gaussian} gives, uniformly on a fixed small
closed real $b$-neighborhood,
\begin{equation}\label{ref:augmented-comparison}
 \begin{split}
 \mathcal A_b(v,\ell)&=(A_bv+C_b\ell,DP_bv),\\
 \|v\|_{\mathcal H^2_a}+|\ell|
   &\le C\|\mathcal A_b(v,\ell)\|_{L^2_a\times\mathbb R^{\dim q}}.
 \end{split}
\end{equation}
The normal charts and the observation support are fixed on every
compact set involved below.

\begin{lemma}[Finite Taylor graphs]\label{ref:finite-taylor}
For every fixed $M\ge1$, let $v^{(M)}(b,x)$ be the coefficientwise
real normal-height polynomial through degree $M$ in the opening
variables, let $\lambda_M$ be the multiplier polynomial, and let $f_M$
be the same truncation of $\mathcal F$.  On a smaller closed real
$b$-neighborhood, each fixed spatial and real parameter derivative of
the normal coefficients is bounded by a fixed polynomial in $1+|Y|$.
The multiplier and action polynomials retain the analytic $b$ dependence
of the prescribed-coordinate jets.  For real parameters,
$\lambda_M-\lambda_b=O(|x|)$.

On every region where the absolute sum of the finite nonconstant height
terms and their first two spatial derivatives is sufficiently small,
the normalized equation of $v^{(M)}$ has an $L^2_a$ remainder at most $C_M|x|^{M+1}$.
When the region contains the observation support, its observation error is
at most $C_M|x|^{M+1}$.  The corresponding
local area integral differs by at most $C_M|x|^{M+1}$ from the integral
of the normal-coordinate density coefficients determined by those
finite normal jets through degree $M$ on that region.  Their complete
integrals sum to $f_M$ for real parameters, and these normal density
coefficients have Gaussian tails bounded by
$C_M(1+r)^{d_M}e^{-c_Mr^2}$ outside $|Y|\le r$.

There is also a globally small real normal graph $T_M$, with height
$\widetilde v_M$ obtained by cutting off $v^{(M)}$ at a radius
$|x|^{-\sigma_M}$, and with $T_M=S_b$
when $x=0$, for which
\begin{align}
 \|\mathcal N_b(\widetilde v_M,\lambda_M)\|_{L^2_a}
       &\le C_M|x|^{M+1},\label{ref:taylor-equation}\\
 |P(T_M)-\mathsf P(p)|+|F(T_M)-f_M(p)|
       &\le C_M|x|^{M+1},\label{ref:taylor-action}\\
 |\nabla_p f_M(p)-J(p)^{\mathsf T}\lambda_M(p)|
       &\le C_M|x|^M,\qquad J=D_p\mathsf P.
       \label{ref:taylor-multiplier}
\end{align}
Its normal height has $C^2$ norm $O(|x|^{1/2})$.  The finitely many
weights $a$ and spatial orders needed in an application are fixed
before $\sigma_M$ is chosen.
\end{lemma}

\begin{proof}
The coefficient equations and their polynomial bounds are those of
Proposition~\ref{lin:family}, including its coefficientwise conversion
to normal graph variables.  The stated absolute smallness keeps every
intermediate finite polynomial at $(b,tx)$, $0\le t\le1$, in the same
normal graph neighborhood.  On such a region, Taylor's formula
for the pointwise geometric graph formula leaves $|x|^{M+1}$ times a
fixed polynomial in $1+|Y|$.  That formula is analytic in height and
first derivatives and affine in the highest second derivatives, so the
local coefficient--module estimates apply to the same derivatives as
in the formal equations.  The compact observation and source maps after
the real graph conversion need only be smooth: their bounded derivatives
through the fixed order $M+1$ give the corresponding finite Taylor
remainders on their fixed support.  Gaussian integration proves the
equation estimate, and the compact Taylor formula proves the observation
estimate.

Expand the normal-coordinate graph area density, including its Gaussian,
in the finite opening variables.  Every fixed coefficient is a
polynomially bounded function times a Gaussian, by the family estimates.
The finite formal change of variables between the prescribed and normal
coordinates identifies their complete coefficient integrals: on an
exhaustion, the change of local density contributes the corresponding
boundary divergence, whose polynomial Gaussian boundary terms vanish
at infinity.  The resulting complete normal-coordinate integrals are
the coefficients of $f_M$.  The local Taylor comparison uses these
normal-coordinate densities and their own tails.  On a real small
normal graph its Taylor remainder is again $|x|^{M+1}$ times a fixed
polynomial and a slightly weaker Gaussian.  This proves the local
area assertion and the coefficient-tail bound.  This assertion concerns
the density coefficients on the complete $S_b$; it does not treat an
uncut polynomial graph as a complete immersed surface.

Choose $D_M$ to bound the growth of the finite height coefficients and
all their derivatives through order two, and any higher fixed orders
used in the application.  Take $\sigma_MD_M<1/2$ and cut off at
$|Y|\asymp |x|^{-\sigma_M}$.  Every nonconstant term and these
derivatives are then $O(|x|^{1/2})$.  The cutoff commutators are supported
at that polynomial radius and have Gaussian norm smaller than every
fixed power of $|x|$.  The interior Taylor estimates just proved give
\eqref{ref:taylor-equation} and \eqref{ref:taylor-action}.
Coefficientwise first variation in \eqref{lin:action-multipliers}
gives \eqref{ref:taylor-multiplier}; an $x$ derivative can lower the
remainder degree by one.  No differentiation of the cutoff with respect
to $x$ is used, and the cut-off graph is not asserted to be analytic in
$x$.
\end{proof}

\begin{lemma}[Localized comparison of real normalized graphs]
\label{ref:real-comparison}
Fix $M$, an admissible weight $a$, and $0<\vartheta<1$, and take
real $(b,x)$ in the stated neighborhood with $|x|$ sufficiently small.
Let $R$ be
large enough that the observation support lies in $|Y|<\vartheta R$,
and let $\chi_R$ equal one there and on $|Y|\le\vartheta R$, vanish
before $|Y|=R$, and satisfy derivative bounds $C_{j,\vartheta}R^{-j}$.
Let $u$ be a smooth real normal graph over $S_b$
on $|Y|<R$, with sufficiently small height and slope.  Let $v$ be
either $v^{(M)}$ in the coefficientwise smallness regime of
Lemma~\ref{ref:finite-taylor} or the height of $T_M$, and require its
$C^2$ norm there to be sufficiently small.  For an arbitrary real
multiplier $\lambda_U$, put
\[
 r_U=\mathcal N_b(u,\lambda_U),\qquad
 e_P=P(u)-\mathsf P(p),\qquad
 w=\chi_R(u-v),\qquad \ell=\lambda_U-\lambda_M,
\]
and define the annular error
\[
 \mathfrak a_R=
 \left\|(1+|Y|)|\nabla(u-v)|+(1+|Y|)^2|u-v|
 \right\|_{L^2_a(\{\vartheta R<|Y|<R\})}.
\]
Then
\begin{equation}\label{ref:real-comparison-estimate}
 \|w\|_{\mathcal H^2_a}+|\ell|
 \le C_{M,\vartheta}\left(\|\chi_R r_U\|_{L^2_a}+|e_P|
                    +|x|^{M+1}+\mathfrak a_R\right).
\end{equation}
The constants are uniform for $b$ in the fixed closed real neighborhood.
In particular, a polynomial bound for the height and its first derivative
on the annulus makes $\mathfrak a_R\le C(1+R)^d e^{-cR^2}$, where one
may take any $c<a\vartheta^2/8$ after absorbing the polynomial factor.

Suppose also that the graph represents the entire part of a real surface
$U$ used inside $|Y|\le\vartheta R$, and that the absolute action of
its omitted part is at most $\tau_{\vartheta R}$.  Then
\begin{equation}\label{ref:real-action-estimate}
 |F(U)-f_M(p)|\le C_M\left(\|w\|_{\mathcal H^2_a}
                         +|x|^{M+1}\right)
                 +\tau_{\vartheta R}
                 +C_{M,\vartheta}(1+R)^{d_M}e^{-c_{M,\vartheta}R^2}.
\end{equation}
For this last conclusion $a$ is chosen so that $a<2c$ for a Gaussian
exponent $c$ retained by the small real graph densities.
\end{lemma}

\begin{proof}
Write $r_M=\mathcal N_b(v,\lambda_M)$.  Before linearization, the exact
subtraction is
\begin{align*}
&\mathcal Q_b(u)-\mathcal Q_b(v)
 +(\mathcal C_b(u)-\mathcal C_b(v))\lambda_b+C_b\ell\\
&\quad=r_U-r_M-(\mathcal C_b(u)-C_b)\ell
 -(\mathcal C_b(u)-\mathcal C_b(v))(\lambda_M-\lambda_b).
\end{align*}
Subtract $A_b(u-v)$ from the first line.  In the quasilinear principal
part retain the coefficients evaluated on $u$ on the second derivatives
of $u-v$.  Their difference from the coefficients of $A_b$ is small
by the height and slope tolerance.  The remaining coefficient differences
multiply the second derivatives of $v$, and their bound is small because
$v$ is $C^2$-small.  The lower terms are controlled in the weighted
slots of \eqref{ref:weighted-H2}.  This arrangement uses no bound for
the second derivatives of $u$.

On the compact source support, $\mathcal C_b(u)-C_b$ is a small bounded
map, so its product with the unknown $\ell$ is absorbed in
\eqref{ref:augmented-comparison}.  The last term in the identity is a
small map on $u-v$, since $\lambda_M-\lambda_b=O(|x|)$.  This argument
does not need an a priori bound for $\lambda_U$.  The observation
equation, with $\chi_R=1$ on its support, gives
\[
 |DP_bw|\le |e_P|+C_M|x|^{M+1}
                 +\varepsilon\|w\|_{\mathcal H^2_a}.
\]
Multiplication by $\chi_R$ creates only height and first-derivative
commutators on its annulus.  The coefficients of the normal graph
equation and its drift have at most the polynomial growth measured by
$\mathfrak a_R$.  Apply \eqref{ref:augmented-comparison}, use
Lemma~\ref{ref:finite-taylor} for $r_M$, and absorb all the small terms.
This proves \eqref{ref:real-comparison-estimate}.

For the action assertion, interpolate the real graphs $v+\tau w$,
$0\le\tau\le1$, on $|Y|<R$; the variation $w$ is compactly supported there.
In the normal tubular
neighborhood, the derivative of their area density obeys
\[
 |D\mathcal L_{v+\tau w}[w]|
 \le C(1+|Y|)^d e^{-c|Y|^2/4}
                         (|w|+|\nabla w|).
\]
The uniform bound for $Y^\perp$ on the real family and the small height
and slope give the retained exponent $c$ and the polynomial prefactor.
Gaussian Cauchy--Schwarz with $a<2c$ bounds its integral by
$C\|w\|_{\mathcal H^2_a}$.  This direct density estimate does not use
stationarity of either graph.  On $|Y|\le\vartheta R$ the interpolated endpoint
is $u$; on the transition annulus all three small graph densities have
Gaussian tails.  Compare the other endpoint with the integrated finite
coefficient densities by Lemma~\ref{ref:finite-taylor}, and add their
full-surface tail and the omitted action $\tau_{\vartheta R}$ of $U$.  The result is
\eqref{ref:real-action-estimate}.
\end{proof}
